% Définition de l'environnement colorboxenv
\NewEnviron{encadre51}[1][]{%
%% Deuxième boite page de garde
\begin{tikzpicture}
    % Node texte (pas de fond ici)
    \node[
        rectangle,
        fill=none,          % <— pas de fond, on va le dessiner derrière
        text=white,
        rounded corners=0pt,
        inner xsep=5pt,
        inner ysep=5pt,
        font=\robotoCondensed
    ] (main) {\fontsize{15}{18}\selectfont#1};


\begin{pgfonlayer}{background}
  \def\padx{0pt}
  \def\pady{0pt}
  \def\r{8pt} % rayon

  % Coins du rectangle extérieur (autour du node)
  \coordinate (SW) at ([xshift=-\padx,yshift=-\pady]main.south west);
  \coordinate (SE) at ([xshift= \padx,yshift=-\pady]main.south east);
  \coordinate (NE) at ([xshift= \padx,yshift= \pady]main.north east);
  \coordinate (NW) at ([xshift=-\padx,yshift= \pady]main.north west);

  \path[fill=primarycolor]
    (SW) --
    % bas : jusqu'à r avant SE
    ([xshift=-\r]SE)
      arc (270:360:\r) % arrondi Bas-Droit, finit sur le bord droit à +r
    --
    % droite : jusqu'à NE
    (NE) --
    % haut : jusqu'à r après NW
    ([xshift=\r]NW)
      arc (90:180:\r)  % arrondi Haut-Gauche, finit sur le bord gauche à -r
    --
    % gauche : retour à SW
    (SW) -- cycle;
\end{pgfonlayer}
\end{tikzpicture}\\
{\itshape\BODY}
}

\NewEnviron{encadre52}[1][]{%%% Deuxième boite page de garde
\begin{tikzpicture}
    % Node texte (pas de fond ici)
    \node[
        rectangle,
        fill=none,          % <— pas de fond, on va le dessiner derrière
        text=white,
        rounded corners=0pt,
        inner xsep=0pt,
        inner ysep=0pt,
        font=\robotoCondensed
    ] (main) {\fontsize{15}{18}\selectfont#1};

    % Fond avec coins "troués" (dessiné DERRIÈRE le texte)
    % Cadre avec UN seul coin arrondi
\begin{pgfonlayer}{background}
  \def\padx{5pt}
  \def\pady{5pt}
  \def\r{8pt} % rayon

  % Coins du rectangle extérieur (autour du node)
  \coordinate (SW) at ([xshift=-\padx,yshift=-\pady]main.south west);
  \coordinate (SE) at ([xshift= \padx,yshift=-\pady]main.south east);
  \coordinate (NE) at ([xshift= \padx,yshift= \pady]main.north east);
  \coordinate (NW) at ([xshift=-\padx,yshift= \pady]main.north west);

  \path[fill=secondarycolor]
    (SW) --
    % bas : jusqu'à r avant SE
    ([xshift=-\r]SE)
      arc (270:360:\r) % arrondi Bas-Droit, finit sur le bord droit à +r
    --
    % droite : jusqu'à NE
    (NE) --
    % haut : jusqu'à r après NW
    ([xshift=\r]NW)
      arc (90:180:\r)  % arrondi Haut-Gauche, finit sur le bord gauche à -r
    --
    % gauche : retour à SW
    (SW) -- cycle;
\end{pgfonlayer}
\end{tikzpicture}\\
{\itshape\BODY}
}
